% !TEX root = thesisPhD.tex 
LaTeX is a document preparation system for high-quality typesetting. It is most often used for medium-to-large technical or scientific documents but it can be used for almost any form of publishing.

\section{An example of text and a list}

LaTeX is based on Donald E. Knuth's TeX typesetting language or certain extensions. LaTeX was first developed in 1985 by Leslie Lamport, and is now being maintained and developed by the LaTeX3 Project. LaTeX is available for free by anonymous ftp \footnote{https://latex-project.org/ftp.html}. LaTeX contains features for:

\begin{itemize}
\item Typesetting journal articles, technical reports, books, and slide presentations.
\item Control over large documents containing sectioning, cross-references, tables and figures.
\item Typesetting of complex mathematical formulas.
\item Advanced typesetting of mathematics with AMS-LaTeX.
\item Automatic generation of bibliographies and indexes.
\item Multi-lingual typesetting.
\item Inclusion of artwork, and process or spot colour.
\item Using PostScript or Metafont fonts.
\end{itemize}

\section{An example of an image inserted}
LaTeX is a powerful typesetting system, used for producing scientific and mathematical documents of high typographic quality. Unlike WYSIWYG tools such as FrameMaker and Word, it uses plain text files that contain formatting
commands. It's big, open source, stable and used by many technical publishing companies. It's also relatively unknown in the technical writing community. This article overviews LaTeX, and directs you to sources of information. Figure \ref{figure:latex-logo} shows the beautiful  logo.

\begin{figure}[htbp]
\begin{center}
\includegraphics[width=3in]{figures/LaTeX_logo}
\caption{The LaTex logo}
\label{figure:latex-logo}
\end{center}
\end{figure}

Avoid using .JPEG, JPG, GIF, BMP or any other bitmap file. Always prefer to use .SVG or .PDF vector files when possible, the result will be much better. After inserting your figure, just continue to writing your text.